\chapter{2008年新课标卷（文）}

\section{单选题}

\begin{problem}
已知集合 \(M = \{x \mid (x + 2)(x - 1) < 0\}\)，\(N = \{x \mid x + 1 < 0\}\)，则 \(M \cap N =\)（\quad）

\choices
    {\((-1, 1)\)}
    {\((-2, 1)\)}
    {\((-2, -1)\)}
    {\((1,2)\)}
\end{problem}

\begin{problem}
双曲线 \(\frac{x^2}{10} - \frac{y^2}{2} = 1\) 的焦距为（\quad）

\choices
    {\(3\sqrt{2}\)}
    {\(4\sqrt{2}\)}
    {\(3\sqrt{3}\)}
    {\(4\sqrt{3}\)}
\end{problem}

\begin{problem}
已知复数 \(z = 1 - \i\)，则 \(\frac{z^2}{z - 1} =\)（\quad）

\choices
    {\(2\)}
    {\(-2\)}
    {\(2\i\)}
    {\(-2\i\)}
\end{problem}

\begin{problem}
设 \( f(x) = x \ln x \)，若 \( f'(x_0) = 2 \)，则 \( x_0 = \)（\quad）

\choices
    {\(\e^2\)}
    {\(\e\)}
    {\(\frac{\ln 2}{2}\)}
    {\(\ln 2\)}
\end{problem}

\begin{problem}
已知平面向量 \(\bs{a} = (1, -3)\)，\(\bs{b} = (4, -2)\)，\(\lambda \bs{a} + \bs{b}\) 与 \(\bs{a}\) 垂直，则 \(\lambda =\)（\quad）

\choices
    {\(-1\)}
    {\(1\)}
    {\(-2\)}
    {\(2\)}
\end{problem}

\begin{problem}
下面的程序框图，如果输入三个实数 \(a\)，\(b\)，\(c\)，要求输出这三个数中最大的数，那么在空白的判断框中，应该填入下面四个选项中的（\quad）

\texfigure[max width=0.65\linewidth]{img_repaint/2008/new_curriculum_liberal/q6_fig1.tex}

\choices
    {\(c > x\)}
    {\(x > c\)}
    {\(c > b\)}
    {\(b > c\)}
\end{problem}

\begin{problem}
已知 \(a_{1} > a_{2} > a_{3} > 0\)，则使得 \((1 - a_i x)^2 < 1\)（\(i = 1,2,3\)）都成立的 \(x\) 取值范围是（\quad）

\choices
    {\(\left(0, \frac{1}{a_1}\right)\)}
    {\(\left(0, \frac{2}{a_1}\right)\)}
    {\(\left(0, \frac{1}{a_3}\right)\)}
    {\(\left(0, \frac{2}{a_3}\right)\)}
\end{problem}

\begin{problem}
设等比数列 \(\{a_{n}\}\) 的公比 \(q = 2\)，前 \(n\) 项和为 \(S_{n}\)，则 \(\frac{S_4}{a_2} =\)（\quad）

\choices
    {\(2\)}
    {\(4\)}
    {\( \frac{15}{2} \)}
    {\(\frac{17}{2}\)}
\end{problem}

\begin{problem}
平面向量 \(\bs{a}\)，\(\bs{b}\) 共线的充要条件是（\quad）

\choices
    {\(\bs{a}, \bs{b}\) 方向相同}
    {\(\bs{a}\)，\(\bs{b}\) 两向量中至少有一个为零向量}
    {\(\exists \lambda \in \R\)，\(\bs{b} = \lambda \bs{a}\)}
    {存在不全为零的实数 \(\lambda_{1}\)，\(\lambda_{2}\)，\(\lambda_{1} \bs{a} + \lambda_{2} \bs{b} = 0\)}
\end{problem}

\begin{problem}
点 \(P(x, y)\) 在直线 \(4x + 3y = 0\) 上，且 \(x\)，\(y\) 满足 \(-14 \leqslant x - y \leqslant 7\)，则点 \(P\) 到坐标原点距离的取值范围是（\quad）

\choices
    {\([0, 5]\)}
    {\([0, 10]\)}
    {\([5, 10]\)}
    {\([5, 15]\)}
\end{problem}

\begin{problem}
函数 \( f(x) = \cos 2x + 2\sin x \) 的最小值和最大值分别为（\quad）

\choices
    {\(-1, 1\)}
    {\(-2, 2\)}
    {\(-3\)，\(\frac{3}{2}\)}
    {\(-2\)，\(\frac{3}{2}\)}
\end{problem}

\begin{problem}
已知平面 \(\alpha \perp\) 平面 \(\beta\)，\(\alpha \cap \beta = l\)，点 \(A \in \alpha\)，\(A \notin l\)，直线 \(AB \parallel l\)，直线 \(AC \perp l\)，直线 \(m \parallel \alpha\)，\(m \parallel \beta\)，则下列四种位置关系中，不一定成立的是（\quad）

\choices
    {\(AB \parallel m\)}
    {\(AC \perp m\)}
    {\(AB \parallel \beta\)}
    {\(AC \perp \beta\)}
\end{problem}

\section{填空题}

\begin{problem}
已知 \( \{a_{n}\} \) 为等差数列，\( a_{3} + a_{8} = 22 \)，\( a_{6} = 7 \)，则 \( a_{5} = \)\fillinblank{}.
\end{problem}

\begin{problem}
一个六棱柱的底面是正六边形，其侧棱垂直底面. 已知该六棱柱的顶点都在同一个球面上，且该六棱柱的高为 \(\sqrt{3}\)，底面周长为 \(3\)，则这个球的体积为\fillinblank{}.
\end{problem}

\begin{problem}
过椭圆 \(\frac{x^2}{5} + \frac{y^2}{4} = 1\) 的右焦点作一条斜率为 \(2\) 的直线与椭圆交于 \(AB\) 两点，\(O\) 为坐标原点，则 \(\triangle OAB\) 的面积为\fillinblank{}.
\end{problem}

\begin{problem}
从甲，乙两品种的棉花中各抽测了 \(25\) 根棉花的纤维长度（单位：\(\mathrm{mm}\)），结果如下：

\begin{center}
\begin{tabular}{c|ccccc}
\hline
品种 & \multicolumn{5}{c}{纤维长度（单位：\(\mathrm{mm}\)）} \\
\hline
甲品种 & \(271\) & \(273\) & \(280\) & \(285\) & \(285\) \\
 & \(287\) & \(292\) & \(294\) & \(295\) & \(301\) \\
 & \(303\) & \(303\) & \(307\) & \(308\) & \(310\) \\
 & \(314\) & \(319\) & \(323\) & \(325\) & \(325\) \\
 & \(328\) & \(331\) & \(334\) & \(337\) & \(352\) \\
\hline
乙品种 & \(284\) & \(292\) & \(295\) & \(304\) & \(306\) \\
 & \(307\) & \(312\) & \(313\) & \(315\) & \(315\) \\
 & \(316\) & \(318\) & \(318\) & \(320\) & \(322\) \\
 & \(322\) & \(324\) & \(327\) & \(329\) & \(331\) \\
 & \(333\) & \(336\) & \(337\) & \(343\) & \(356\) \\
\hline
\end{tabular}
\end{center}

由以上数据设计了如下茎叶图

\begin{center}
\begin{tabular}{rrrrr|c|lllllll}
\multicolumn{5}{c|}{甲} & & \multicolumn{7}{c}{乙} \\
\hline
 &  &  & \(3\) & \(1\) & \(27\) &  &  &  &  &  &  &  \\
 & \(7\) & \(5\) & \(5\) & \(0\) & \(28\) & \(4\) &  &  &  &  &  &  \\
 &  & \(5\) & \(4\) & \(2\) & \(29\) & \(2\) & \(5\) &  &  &  &  &  \\
\(8\) & \(7\) & \(3\) & \(3\) & \(1\) & \(30\) & \(4\) & \(6\) & \(7\) &  &  &  &  \\
 &  & \(9\) & \(4\) & \(0\) & \(31\) & \(2\) & \(3\) & \(5\) & \(5\) & \(6\) & \(8\) & \(8\) \\
 & \(8\) & \(5\) & \(5\) & \(3\) & \(32\) & \(0\) & \(2\) & \(2\) & \(4\) & \(7\) & \(9\) &  \\
 &  & \(7\) & \(4\) & \(1\) & \(33\) & \(1\) & \(3\) & \(6\) & \(7\) &  &  &  \\
 &  &  &  &  & \(34\) & \(3\) &  &  &  &  &  &  \\
 &  &  &  & \(2\) & \(35\) & \(6\) &  &  &  &  &  &  \\
\end{tabular}
\end{center}

根据以上茎叶图，对甲，乙两品种棉花的纤维长度作比较，写出两个统计结论：

\begin{circlelist}
\item\fillinblank{}；
\item\fillinblank{}.
\end{circlelist}
\end{problem}

\begin{problem}
如图，\(\triangle ACD\) 是等边三角形，\(\triangle ABC\) 是等腰直角三角形，\(\angle ACB = 90^{\circ}\)，\(BD\) 交 \(AC\) 于 \(E\)，\(AB = 2\).

\begin{enumerate}
\item 求 \(\cos \angle CBE\) 的值；
\item 求 \(AE\).
\end{enumerate}

\bitmapfigure[max width=0.34\linewidth]{img/2008/new_curriculum_liberal/q17_fig1.png}

\end{problem}

\begin{problem}
如下的三个图中，上面的是一个长方体截去一个角所得多面体的直观图. 它的正视图和俯视图在下面画出（单位：\(\mathrm{cm}\)）.

\begin{enumerate}
\item 在正视图下面，按照三视图的要求画出该多面体的俯视图；
\item 按照给出的尺寸，求该多面体的体积；
\item 在所给直观图中连结 \(BC'\)，证明：\(BC' \parallel\) 面 \(EFG\).
\end{enumerate}

\bitmapfigure[max width=0.72\linewidth]{img/2008/new_curriculum_liberal/q18_fig1.png}

\end{problem}

\begin{problem}
为了了解《中华人民共和国道路交通安全法》在学生中的普及情况，调查部门对某校\(6\)名学生进行问卷调查. \(6\)人得分情况如下：\(5,6,7,8,9,10\). 把这\(6\)名学生的得分看成一个总体.

\begin{enumerate}
\item 求该总体的平均数；
\item 用简单随机抽样方法从这 \(6\) 名学生中抽取 \(2\) 名，他们的得分组成一个样本. 求该样本平均数与总体平均数之差的绝对值不超过 \(0.5\) 的概率.
\end{enumerate}
\end{problem}

\begin{problem}
设函数 \( f(x) = ax - \frac{b}{x} \)，曲线 \( y = f(x) \) 在点 \((2, f(2))\) 处的切线方程为 \( 7x - 4y - 12 = 0 \).

\begin{enumerate}
\item 求 \( f(x) \) 的解析式；
\item 证明：曲线 \( y = f(x) \) 上任一点处的切线与直线 \( x = 0 \) 和直线 \( y = x \) 所围成的三角形面积为定值，并求此定值.
\end{enumerate}
\end{problem}

\begin{problem}
如图，过圆 \(O\) 外一点 \(M\) 作它的一条切线. 切点为 \(A\)，过 \(A\) 点作直线 \(AP\) 垂直直线 \(OM\)，垂足为 \(P\).

\begin{enumerate}
\item 证明：\(OM \cdot OP = OA^2\)；
\item \(N\) 为线段 \(AP\) 上一点，直线 \(NB\) 垂直直线 \(ON\)，且交圆 \(O\) 于 \(B\) 点. 过 \(B\) 点的切线交直线 \(ON\) 于 \(K\). 证明：\(\angle OKM = 90^{\circ}\).
\end{enumerate}

\bitmapfigure[max width=0.48\linewidth]{img/2008/new_curriculum_liberal/q22_fig1.png}

\end{problem}

\begin{problem}
已知 \(m \in \R\)，直线 \(l: mx - (m^2 + 1)y = 4m\) 和圆 \(C: x^2 + y^2 - 8x + 4y + 16 = 0\).

\begin{enumerate}
\item 求直线 \(l\) 斜率的取值范围；
\item 直线 \(l\) 能否将圆 \(C\) 分割成弧长的比值为 \(\frac{1}{2}\) 的两段圆弧？为什么？
\end{enumerate}
\end{problem}

\begin{problem}
已知曲线 \(C_{1}\)： \(\left\{ \begin{array}{l}x = \cos \theta ,\\ y = \sin \theta . \end{array} \right.\) （ \(\theta\) 为参数），曲线 \(C_2\)： \(\left\{ \begin{array}{l}x = \frac{\sqrt{2}}{2} t - \sqrt{2},\\ y = \frac{\sqrt{2}}{2} t, \end{array} \right.\) （ \(t\) 为参数）.

\begin{enumerate}
\item 指出 \(C_1\)，\(C_2\) 各是什么曲线，并说明 \(C_1\) 与 \(C_2\) 公共点的个数；
\item 若把 \(C_1\)，\(C_2\) 上各点的纵坐标都压缩为原来的一半，分别得到曲线 \(C_1'\)，\(C_2'\). 写出 \(C_1'\)，\(C_2'\) 的参数方程. \(C_1'\) 与 \(C_2'\) 公共点的个数和 \(C_1\) 与 \(C_2\) 公共点的个数是否相同？说明你的理由.
\end{enumerate}
\end{problem}
